\documentclass[10pt,a4paper]{amsart}
\usepackage[margin=19mm]{geometry}
\usepackage{amsmath,amssymb,mathtools}
\usepackage[scr=rsfs,cal=euler]{mathalfa}
\usepackage{xfrac}
\usepackage[unicode,hidelinks,bookmarks=false]{hyperref}
\newcommand{\ie}{i.e.\ }
\newcommand{\cf}{cf.\ }
\newcommand{\resp}{respectively\ }
\newcommand{\Subsection}{\subsection}
\providecommand{\nobreakdash}{\nobreak\hbox{-}}
\setlength{\emergencystretch}{2em}
\multlinegap0pt
\usepackage{bm}
\usepackage{bbm}
\usepackage{dsfont}
\usepackage{xspace}
\usepackage{cancel}
\usepackage{mathtools}
\usepackage{amsthm}
\makeatletter
\@ifundefined{theorem}{\newtheorem{theorem}{Theorem}}{}
\@ifundefined{corollary}{\newtheorem{corollary}[theorem]{Corollary}}{}
\makeatother
\makeatletter
\newcommand{\R}{\mathbb R}
\newcommand{\bq}{\begin{equation}}
\newcommand{\calP}{\mathcal P}
\newcommand{\diag}{{\rm diag}}
\newcommand{\dy}{\textnormal{d}y}
\newcommand{\eq}{\end{equation}}
\renewcommand{\ge}{\geqslant}
\renewcommand{\geq}{\geqslant}
\newcommand{\intr}{\int_{\R^d}}
\renewcommand{\le}{\leqslant}
\newcommand{\lt}{\left}
\newcommand{\nabA}{\nabla_{\!\!{}_A}}
\newcommand{\nabS}{\nabla_{\!\!{}_S}}
\newcommand{\pa}{\partial}
\expandafter\ifx\csname remark\endcsname\relax
\newenvironment{remark}{{\flushleft {\bf Remark.}}}{}
\fi
\expandafter\ifx\csname remarks\endcsname\relax
\newenvironment{remarks}{{\flushleft {\bf Remarks.}}}{}%\newtheorem{definition}{Definition}[section]
\fi
\newcommand{\rt}{\right}
\newcommand{\supp}{{\rm supp}}
\makeatother
\title[Finite-time breakdown of the Euler-alignment system for supe]{Finite-time breakdown of the Euler-alignment system for supercritical initial data}
\author{Young-Pil Choi, Eitan Tadmor}
\date{Source: arXiv:2607.08060v1 (CC BY 4.0)}
\begin{document}
\maketitle

\noindent\emph{Source and licence.} Finite-time breakdown of the Euler-alignment system for supercritical initial data. Version arXiv:2607.08060v1 (CC BY 4.0), distributed under the Creative Commons Attribution 4.0 International licence, as recorded in the arXiv licence metadata for this exact version and shown on the abstract page https://arxiv.org/abs/2607.08060v1. Full rights notice: licenses/arxiv-2607.08060-CC-BY-4.0.txt.\par\smallskip

\noindent\textit{Editorial scope.} Counted unit: the Corollary whose source label is \texttt{cor:gen\_sup\_const} in the article (source0.tex lines 1095--1125, source label \texttt{cor:gen\_sup\_const}), delivered together with the article's own complete original proof of it (source0.tex lines 1146--1220, one \texttt{proof} environment of 75 lines). No mathematics is written, repaired or re-derived: statement and proof are the article's own text line for line. The proof is detached from the statement in the source: the article states the result at lines 1095--1125 and prints its proof at lines 1146--1220, and the proof environment's own header names the statement, so the association is the article's own. Required context retained uncounted: LA:const (lines 579--584). Every definition, display and label of those lines is retained unchanged. Omitted from this package and not redistributed: 1--578, 585--1094, 1126--1145, 1221--1754. Nothing inside a retained excerpt is omitted or altered: every retained body is delivered exactly as the article's lines are written, so no omission receipt of any kind accompanies this package. Citations: the retained excerpts carry no citation command at all, so no bibliography entry of the article is transcribed and no reference is redistributed. Reference adaptation: none was needed and none was made; every reference inside the delivered unit resolves inside it, so no unresolved reference or citation is produced. Printed slips kept literal: no printed word, symbol, hypothesis or number of the counted statement or of its proof was altered. Layout and class: the article's own class is \texttt{amsart} with packages including amsmath, amssymb, amsfonts, mathtools, amsthm, enumitem, hyperref, geometry, mathrsfs, graphicx, sidecap, rotating, dsfont, ulem; the wrapper is the lane's pinned \texttt{amsart} editorial wrapper at 10pt on A4, which restates the theorem-like environments the retained text needs and adds the article's own macro definitions that the retained text needs (\texttt{\textbackslash R}, \texttt{\textbackslash bq}, \texttt{\textbackslash calP}, \texttt{\textbackslash diag}, \texttt{\textbackslash dy}, \texttt{\textbackslash eq}, \texttt{\textbackslash ge}, \texttt{\textbackslash geq}, \texttt{\textbackslash intr}, \texttt{\textbackslash le}, \texttt{\textbackslash lt}, \texttt{\textbackslash nabA}, \texttt{\textbackslash nabS}, \texttt{\textbackslash pa}, \texttt{\textbackslash rt}, \texttt{\textbackslash supp}), with extra wrapper packages (bm, bbm, dsfont, xspace, cancel, mathtools). The delivered numbering is the unit's own. Assets: no retained span contains a figure environment, a graphics-inclusion command, a PDF-page inclusion, a table or a TikZ picture; no image file is copied into this package, the package declares no asset dependency and no image file is redistributed with it. Licence: version arXiv:2607.08060v1 of this article is distributed under the Creative Commons Attribution 4.0 International licence, as recorded in the arXiv licence metadata for this exact version and shown on the abstract page https://arxiv.org/abs/2607.08060v1. A full rights notice is delivered at licenses/arxiv-2607.08060-CC-BY-4.0.txt. Characters: every retained excerpt is pure ASCII, so the delivered engine, XeLaTeX with its default Latin Modern text font, reports no missing character.
 \begin{align}\label{LA:const}
\begin{aligned}
&\pa_t \eta(t,x) = v(t,x), \quad t>0, \ x \in \supp(\rho_0), \cr
&\pa_t v(t,x)  = \kappa \intr   (v(t,y) - v(t,x))\,\rho_0(\dy),
\end{aligned}
\end{align}

\begin{corollary} \label{cor:gen_sup_const}
Let $d \geq 3$ and consider the system \eqref{LA:const}. Assume $\rho_0\in C^0\cap\calP(\R^d)$ and $u_0\in C^1\cap W^{1,\infty}(\R^d;\R^d)$. Suppose that there exist $x_*\in\supp(\rho_0)$, $\alpha_1\in(0,\frac1\kappa)$, and $\delta>0$ such that, setting
\[
M_*:=\nabla u_0(x_*),\quad S_*:=\nabS u_0(x_*),\quad A_*:=\nabA u_0(x_*),
\]
and denoting by $\mu_1\le\cdots\le\mu_d$ the eigenvalues of $S_*$, one has
\bq\label{eq:sep0}
|1+\alpha_1\mu_j|\ge \delta \quad\text{for all }j=1,\dots,d,
\eq
and the number
\[
r:=\#\{j\in\{1,\dots,d\}:1+\alpha_1\mu_j<0\}
\]
is odd. Assume moreover the smallness condition
\bq\label{eq:sma0}
\alpha_1\|A_*\|<\delta.
\eq
Then there exists $\alpha_*\in(0,\alpha_1)$ such that
\[
\det(I_d+\alpha_*M_*)=0.
\]
Consequently, there exists
\[
t_*\in\lt(0,-\frac1\kappa\log(1-\kappa\alpha_1)\rt)
\]
such that
\[
\det\nabla\eta(t_*,x_*)=0.
\]
In particular, the Lagrangian flow degenerates at $x_*$ in finite time.
\end{corollary}

\begin{proof}[Proof of Corollary \ref{cor:gen_sup_const}]
We proceed in several steps. 

\medskip
\noindent
\textit{Step 1: Reduction to an algebraic root in $\alpha$.}
We first recall
\[
\nabla\eta(t,x_*)=I_d+\alpha(t)M_*, \quad \alpha(t)=\frac{1-e^{-\kappa t}}{\kappa}\in \lt[0,\frac1\kappa\rt).
\]
Let $f(\alpha):=\det(I_d+\alpha M_*)$ for $\alpha\in[0,1/\kappa)$. Then $\det\nabla\eta(t,x_*)=f(\alpha(t))$, and $f$ is continuous with $f(0)=1>0$. Hence it suffices to prove $f(\alpha_1)<0$, since then the intermediate value theorem yields $\alpha_*\in(0,\alpha_1)$ with $f(\alpha_*)=0$, and the monotonicity of $\alpha(t)$ provides $t_*>0$ such that $\alpha(t_*)=\alpha_*$.

\medskip
\noindent
\textit{Step 2: Diagonalization of $S_*$ and sign of the unperturbed determinant.}
Since $S_*$ is symmetric, there exists an orthogonal matrix $Q\in O(d)$ such that
\[
Q^\top S_*Q=\Lambda:=\diag(\mu_1,\dots,\mu_d).
\]
Set $\widetilde A_*:=Q^\top A_*Q$, which remains skew-symmetric, and note that $\|\widetilde A_*\|=\|A_*\|$.
By invariance of the determinant under orthogonal similarity transformations,
\[
f(\alpha_1)=\det \lt(I_d+\alpha_1(\Lambda+\widetilde A_*)\rt).
\]
Define the diagonal matrix
\[
D:=I_d+\alpha_1\Lambda=\diag(d_1,\dots,d_d), \quad d_j:=1+\alpha_1\mu_j,
\]
so that
\[
f(\alpha_1)=\det(D+\alpha_1\widetilde A_*).
\]
By the separation condition \eqref{eq:sep0}, we have $d_j\neq0$ for all $j$ and hence $D$ is invertible.
Moreover, since exactly $r$ of the $d_j$ are negative and $r$ is odd, it follows that
\[
\det D=\prod_{j=1}^d d_j= -\prod_{j=1}^d |d_j|< 0.
\]

\medskip
\noindent
\textit{Step 3: Stability of the sign under a small skew perturbation.}
We write
\[
f(\alpha_1)=\det(D+\alpha_1\widetilde A_*)
=\det D \cdot \det(I_d+X),
\quad
X:=D^{-1}(\alpha_1\widetilde A_*).
\]
Since $\|D^{-1}\|\le 1/\delta$, we obtain
\[
\|X\|\le \frac{\alpha_1\|A_*\|}{\delta}<1
\]

This implies that every eigenvalue $\lambda$ of $X$ satisfies $|\lambda|<1$. Hence no real eigenvalue can cross $-1$, and complex eigenvalues occur in conjugate pairs. Therefore
\[
\det(I_d+X)=\prod_{\lambda\in\sigma(X)}(1+\lambda)>0.
\]
It follows that
\[
f(\alpha_1)=\det D \cdot \det(I_d+X)<0.
\]

\medskip
\noindent
\textit{Step 4: Conclusion.}
Since $f(0)=1>0$ and $f(\alpha_1)<0$, there exists $\alpha_*\in(0,\alpha_1)$ such that
\[
f(\alpha_*)=0.
\]
Finally, by the strict monotonicity of $\alpha(t)$, there exists $t_*>0$ such that $\alpha(t_*)=\alpha_*$, and hence
\[
\det\nabla\eta(t_*,x_*)=0.
\]
This completes the proof.
\end{proof}

\end{document}
